\begin{center}
    {\large \textbf{Giorno 7} \textbar\ Betel Quid \textbar\ 15 Agosto 2024 \textbar\ \textbf{Kuta}, Lombok}
\end{center}

\noindent\rule{\textwidth}{0.4pt}
\begin{figure}[h!] % Portrait images below
    \centering
    % First portrait image (on the left)
    \begin{minipage}[h]{0.48\textwidth} % Set width equal to half the page
        \centering
        \includegraphics[width=\textwidth]{images/flags/Screenshot villaggi Lombok.png}
    \end{minipage}
    \hfill
    % Text
    \begin{minipage}[h]{0.48\textwidth}
        \raggedright
        \textbf{Superficie}: 4.567 km² \\
        \vspace{0.3cm}
        \textbf{Popolazione}: 3,3 milioni \\
        \vspace{0.3cm}
        \textit{L'isola}
    \end{minipage}
\end{figure}

\landscapeLargeMargin{images/day7/PXL_20240815_001756799.jpg}

\landscapeLargeMargin{images/day7/PXL_20240815_020838116.MP.jpg}

\landscapeLargeMargin{images/day7/PXL_20240815_020928473.jpg}

\landscapeLargeMargin{images/day7/PXL_20240815_023746808.jpg}

\landscapeLargeMargin{images/day7/PXL_20240815_030131561.MP.jpg}

\landscapeLargeMargin{images/day7/PXL_20240815_030135418.MP.jpg}

\noindent 15 Agosto 2024 \newline

Se ne son dette tante su questo Novotel, ma come si dorme qui non si dorme da nessuna parte a Lombok. Tutta la notte di fila, senza animali o preghiere a svegliarci, ci sentiamo rigenerati. Anche la colazione merita una menzione speciale, pur trattandosi di un buffet occidentale. Gustiamo tutto quello che hanno da offrire, stando lontano dai cibi piccanti. Il caffè è prodotto localmente ed è ottimo, anche se non è fatto miscelando acqua calda con caffè macinato come fanno di solito i locali, ma è filtrato all’americana. Un po' rimpiango il metodo tradizionale indonesiano.

Dopo esserci abbuffati di pancake, paghiamo la colazione con una trasfusione e torniamo in camera a prepararci. No, in realtà non è vero, non è per niente caro ma rispetto a quello che abbiamo visto nei villaggi nativi nei giorni precedenti qui i prezzi sono decisamente più alti. Torniamo in camera a prepararci e, conoscendo la puntualità degli indonesiani, ci rechiamo all'appuntamento all'ingresso dell'hotel con qualche minuto di anticipo.

Mentre attendiamo la nostra macchina chiacchieriamo con un dipendente del Novotel che parla italiano e quando gli descriviamo il tour dei tre villaggi che abbiamo organizzato, ci fa notare che due villaggi sono vicinissimi tra loro e di solito vengono visitati in alternativa, mai assieme. Iniziamo a pensare di essere caduti in uno scam, visto che i tour normali raggiungono anche un villaggio dove si lavora la ceramica, più lontano e quindi più dispendioso per loro.

Alle 9:10 l'autista ancora non si vede e chiamo il suo capo, già pronto a smadonnare. Mi assicura che arriverà a breve e così è per fortuna. Il nostro autista si chiama Abi e scopriamo che è in realtà il proprietario del centro turistico. Inizio a interrogarlo: dapprima per cercare di stanare le fregature — quindi quali villaggi vediamo? Fino a che ora stiamo? Poi ci porti qui, giusto? E stiamo fino al tramonto, vero? — poi mi faccio dare qualche consiglio sulle isole Gili, dal momento che lui ha fatto lì la guida turistica per cinque anni prima di spostarsi a Kuta. Ci racconta che le tre isole hanno ciascuna il suo soprannome: Gili Trawangan è la Party Island, Gili Meno è la Honeymoon Island e Gili Air la Relax Island. Gli chiediamo delle tartarughe e ci dice di rivolgerci alle guide più anziane ed esperte, perché i più giovani non danno garanzia di poterle vedere, e ci consiglia di partire col tour dalla Gili Air "next to Alibaba restaurant". Per evitare possibili scam lo incentivo promettendogli di portargli clienti, dal momento che Enrico e Cristina verranno a Kuta e avranno bisogno di guide e trasporti.

Il primo villaggio nel quale ci porta si chiama Sukarara. All'ingresso siamo accolti da una guida locale, il caro Holidin, che ci propone un tour all'interno. Scopriamo subito che, a differenza del Tetebatu, queste zone sono molto più aride: non possono godere dell'acqua del lago del Monte Rinjani portata a valle, pertanto il riso cresce più lentamente e hanno bisogno di strutture apposite per conservarlo durante l'anno. Riescono a raccoglierlo solo una volta l'anno, a differenza dei villaggi più a nord che arrivano a due o tre coltivazioni grazie alle fonti d'acqua.

Dapprima scattiamo qualche foto a una ricostruzione delle abitazioni tipiche ed entriamo a esplorarla. Ci mostrano poi un Lumbung, una casa adibita alla conservazione del riso, condivisa tipicamente da cinque famiglie così da mettere in comune le scorte. Scopriamo finalmente anche il nome delle piattaforme di bambù che abbiamo visto ovunque fin dall'inizio: si chiamano Berugak, mentre sull'isola di Java le chiamano semplicemente Gazebo. Sono realizzate con due specie di bambù locali: il green bambù, più solido, per la struttura, e il geko bambù, più morbido, per i tetti.

Prima ancora di entrare nel villaggio veniamo intercettati da un'altra guida, Rudy, che conosce qualche parola in italiano e viene a fare pratica con noi. È tuttavia molto rispettoso e ci lascia nelle mani di Holidin, che ci aveva accolti per primo. Holidin è ferratissimo sui nomi delle piante, ma per fortuna ci siamo allenati nel Tetebatu e incredibilmente riconosco subito il chili e il dragon fruit, con sua grande sorpresa. Dopo non ne azzeccherò più neanche una ovviamente. Quando arriviamo alla pianta di cacao — che sarà Elisa a riconoscere — ci dice che loro non la usano come colorante per i tessuti ma la mangiano perché ne sono molto ghiotti.

Assistiamo alla macinatura del caffè: anche qui lo macinano insieme al riso, lo chiamano Lombok coffee e non lo filtrano mai, ma lo mescolano direttamente con l'acqua calda. Ci consiglia di provarlo con lo zucchero di canna naturale. Agli abitanti di Lombok il caffè piace molto zuccherato e questo è un problema: amano i dolci e tutti i loro prodotti abbondano di zucchero, così fin da piccoli i bambini rischiano di sviluppare il diabete.

Vediamo piante di tamarindo, papaya, mango — di cui esistono due specie — e in particolare l'honey mango, chiamato così perché è particolarmente dolce. Nelle case tradizionali di bambù e legno, le donne tessono i Sarong per otto ore al giorno. Assistiamo alla tessitura con i loro strumenti originali e questa volta fanno provare anche Elisa. Ci spiegano le tecniche per colorare il tessuto prima di lavorarlo e notiamo con sorpresa che qui utilizzano alimenti diversi rispetto al Tetebatu: per il verde gli spinaci, per l'arancio la curcuma, per il nero il caffè, per il rosso il chili, per il rosa il dragon fruit, per il marrone il tamarindo, e solo per il grigio e i colori più difficili usano tessuti sintetici.

Le donne si staccano dalla tessitura solo per preparare da mangiare alla famiglia. Le bambine iniziano a tessere da piccole, intorno ai dieci o dodici anni, e impiegano almeno tre anni per diventare brave. Secondo la tradizione, una donna per sposarsi deve essere in grado di tessere un Sarong da sola — non è qualcosa che si impara a scuola, ma viene tramandata di generazione in generazione. Le bambine lo fanno dopo scuola ogni giorno, iniziando a mezzogiorno e continuando fino a sera, indossando spesso fino a tre cinture posturali perché è un'attività faticosa che fa male alla schiena. Lo fanno sette giorni su sette, senza vacanze o weekend, fermandosi solo per pregare o per preparare il cibo ai mariti.

I bambini vanno a scuola tipicamente fino alle superiori, a diciotto anni, e solo i figli delle famiglie più ricche accedono all'università. A Lombok ce ne sono tre, tutte nella capitale Mataram. Chi ha più possibilità si sposta sulle altre isole, come a Jakarta o Yogyakarta su Java.

Una delle ultime sorprese all'interno del villaggio ci lascia senza parole: appesi agli alberi ci sono dei tronchi cavi pieni, pieni, pieni di api. Qui le tengono sugli alberi invece che nelle arnie, e a differenza delle nostre questa specie non punge.

Prima di uscire dal villaggio vediamo il secondo tipo di mango, il manalaghi, che in indonesiano significa "Where is more?" — perché è molto buono e le persone ne vogliono sempre di più.

Terminato il tour visitiamo il centro turistico dove ci mostrano altri Sarong e continuano a riempirci di nozioni sulla tessitura. Una decorazione molto tradizionale e difficile da realizzare rappresenta due pavoni che si guardano l'uno con l'altro: è il simbolo della luna di miele per questo villaggio. Le donne devono realizzare tre Sarong decorati con i pavoni per potersi sposare — uno per il marito, uno per la suocera e uno per la sposa stessa. Considerando che ci vogliono almeno due settimane a tesserne uno, per prepararsi al matrimonio serve almeno un mese e mezzo.

Siamo accompagnati poi a vedere i vestiti tradizionali del matrimonio e ce li fanno indossare per una foto ricordo. Si spiega così perché durante tutto il tour ci fosse una donna rimasta in silenzio fino a quel momento: siccome ero con Elisa, serviva una donna per vestire un'altra donna. Il mio vestito comprende anche un pugnale caratteristico che secondo la tradizione viene usato per proteggere la sposa: si chiama Pumajù, tipico di Lombok, col manico realizzato in tamarindo. Il pugnale tipico indonesiano si chiama invece Karis.

Terminate le foto e i ringraziamenti è il momento di salutare la nostra guida, che ci congeda con un messaggio che merita di essere riportato: "If you want you can leave us some tips, but if you don't want don't worry: we do it because we love tourist." Non esiste frase migliore per guadagnarsi una mancia.

Mentre scattiamo qualche foto all'uscita veniamo intercettati nuovamente da Rudy, la guida che parla un po' di italiano. Ci intrattiene con alcune frasi imparate a memoria e, dal momento che oggi è il 15 agosto, gli suggerisco di accogliere i turisti italiani dicendo "Buon Ferragosto", spiegandogli che da noi oggi è festa nazionale. Probabilmente colpito da questa semplice gentilezza mi chiede se può farmi alcune domande e inizia a interrogarmi su come si dicano in italiano alcune frasi standard che è solito dire in inglese ai turisti. Traduco pazientemente, lui sembra memorizzare ogni cosa e ripete ogni frase con grande attenzione. Infine mi chiede il numero di telefono, nel caso avesse bisogno di aiuto con altre frasi. Ovviamente acconsento: la situazione mi sembra surreale, ma l'estrema gentilezza dei suoi modi mi fa venire voglia di essergli d'aiuto.

Queste persone sono tutte estremamente gentili con noi e, come ha detto la nostra guida, si vede che non lo fanno per soldi in primo luogo, ma perché sono stati educati così. Dal canto nostro cerchiamo di ricambiare: li salutiamo ogni volta, chiediamo come stanno, cerchiamo di ricordarci il loro nome e facciamo attenzione alle loro usanze e alla loro cultura.

Non siamo tutti così purtroppo. A un certo punto arriva un pullman carico di turisti italiani. Ne scende una coppia, che subito raggiunge una donna intenta a tessere. La ragazza si allontana subito disinteressata. mentre il compagno tira fuori il cellulare, lo punta in faccia alla donna da meno di un metro di distanza ed esclama "Faccio un video a questa e arrivo." Questa scena mi fa venire in mente la parabola di Tatanaele e Isappo:

\FloatBarrier
\landscapeLargeMargin{images/day7/tatanaele-e-isappo.jpg}
\FloatBarrier

\landscapeLargeMargin{images/day7/PXL_20240815_042530442.jpg}

\landscapeLargeMargin{images/day7/IMG_4756.JPG}

\landscapeLargeMargin{images/day7/PXL_20240815_052945302.jpg}

\landscapeLargeMargin{images/day7/PXL_20240815_053435142.jpg}

\landscapeLargeMargin{images/day7/PXL_20240815_053741318.jpg}

\landscapeLargeMargin{images/day7/PXL_20240815_060352994.jpg}

\landscapeLargeMargin{images/day7/PXL_20240815_061957390.jpg}

\landscapeLargeMargin{images/day7/PXL_20240815_064724517.jpg}

\landscapeLargeMargin{images/day7/PXL_20240815_070108612.jpg}

\landscapeLargeMargin{images/day7/PXL_20240815_064808722.MP.jpg}

\landscapeLargeMargin{images/day7/PXL_20240815_070043162.MP.jpg}

\landscapeLargeMargin{images/day7/PXL_20240815_072753846.jpg}

\landscapeLargeMargin{images/day7/PXL_20240815_072835961.MP.jpg}

\landscapeLargeMargin{images/day7/PXL_20240815_073528559.jpg}

Ecco: quel turista italiano meriterebbe di prenderle come Tatanaele.

\begin{center}
    % Decorative line
    \noindent\rule{0.6\textwidth}{0.4pt}
\end{center}

Usciamo dal centro turistico e ci aggiriamo per il mercato locale lungo la strada principale. Pare che si tenga ogni giovedì e quindi siamo fortunati a potervi assistere — chissà se sarà vero o se ce lo dicono per invogliarci a spendere. In ogni caso giriamo per le bancarelle scattando foto a quelle che vendono soprattutto cibo, frutta e artigianato, e rientro in modalità bodyguard mentre Elisa scatta le foto. Per fortuna oggi non sembro un surfista albino: occhiali scuri, maglietta verde militare, zaino tattico nero, pantaloncini beige con tasconi e scarpe da trekking. Mi ricorda il mio outfit in Messico e mi sento più a mio agio.

Tra una bancarella e l'altra passiamo davanti a una scuola, dove i bambini che giocano fuori ci chiedono "Money! Money!" Ormai stanno smantellando le bancarelle quindi torniamo alla macchina e contrattiamo con Abi un cambio di programma: la spiaggia che dovremmo vedere a fine giornata ci sembra troppo turistica e dalle recensioni online pare anche piuttosto sporca, così gli chiediamo se per lo stesso prezzo può portarci a vedere il pottery village, il villaggio dove lavorano la terracotta. Accetta subito — probabilmente perché capisce che così avrebbe terminato prima la giornata — tuttavia, forse per incomprensione o più probabilmente perché aveva qualche provvigione coi produttori locali, ci porta in un magazzino di prodotti di terracotta dove non vediamo la lavorazione ma possiamo solo fare acquisti.

Poco male, in realtà, perché volevamo portarci a casa un souvenir e questo è probabilmente il posto più adatto dove trovare buoni prezzi prima delle isole Gili e di Bali. Il magazzino si trova in un capannone enorme e ospita sia cianfrusaglie — i classici souvenir con la scritta "Lombok" — che prodotti di terracotta di pregiata rifinitura. Dopo aver girato e contrattato un po' ci portiamo a casa un servizio composto da teiera, vassoio e quattro tazzine per il caffè di Lombok, e un vaso realizzato con terracotta e gusci d'uovo. Il prezzo è contenuto e dopotutto almeno un vaso l'avevo messo in conto. La sera obbligherò poi Elisa a giurarmi: "basta vasi." Ci imballano il tutto in uno scatolone coi giornali e già temo il momento in cui dovremo imbarcarlo sull'aereo.

\begin{center}
    % Decorative line
    \noindent\rule{0.6\textwidth}{0.4pt}
\end{center}

Salutiamo e ci dirigiamo verso la nuova destinazione: l'antico villaggio di Ende, pronunciato Nde. Abi ci ha fatto capire che con un tour privato l'autista diventa la tua guida personale, così attacco bottone e mi racconta qualche fatto curioso: Lombok ha l'85% del territorio coperto da campi di riso; oppure, se bevi il caffè di Lombok per una settimana e smetti la settimana successiva ti viene il mal di testa. Vedremo quando saremo tornati. Gli chiedo anche se è possibile trovare un pugnale come quello provato col vestito tradizionale e mi dice che ci pensa lui. Non ne saprò più niente.

Durante il viaggio vengo contattato da Rudy che mi chiede di tradurgli alcune frasi in italiano per i turisti. Nei messaggi usa vari appellativi — "sir", "bro", "friend" — e chiude con un "Thank you very much Mr Marco."

Arriviamo al villaggio di Ende e anche qui siamo accolti da una guida, Amat, un abitante del villaggio che ci propone un tour. Accettiamo volentieri, reduci dalla positiva esperienza con Holidin. A differenza di prima, qui il tour ha un prezzo vero e proprio — che coincide però con la mancia che avevo lasciato al suo predecessore, quindi poco male.

La prima caratteristica che mi colpisce delle loro case è quanto i tetti siano bassi: la guida ci spiega che un ospite deve entrare con rispetto, quindi inchinandosi. Io rischio una testata e Amat mi soprannomina Buffon. La casa viene chiamata Balitani, dall'unione di Bali — Casa — e Tani — Famiglia. All'interno ci sono due stanze: una con pareti e senza finestre, dove dormono le donne, e una esterna coperta solo dal tetto, dove dormono gli uomini. La tradizione vuole che gli sposi dormano insieme all'interno solo durante la luna di miele, per poi separarsi di nuovo fino a quando non vogliono avere un altro figlio. Il suolo è realizzato con argilla e sterco di vacca, quindi le case più nuove hanno un odore particolare mentre quelle più antiche no. Non hanno letti o materassi, ma dormono su stuoie con piccoli cuscini. Il tetto è fatto di erba — ottimo per far scorrere la pioggia, ma va cambiato ogni sei o sette anni. All'interno c'è anche una cucina dove le donne cucinano solo con il fuoco, lentamente, per non far prendere fuoco al tetto. Ogni casa è per una famiglia sola; quando i figli si sposano viene costruita una nuova casa per i novelli sposi.

Ci dice che le villette del Novotel sono strutturate come la sua casa e ci chiede se per caso dormiamo nella stanza 505. Eh no, my friend.

Anche qui troviamo le strutture per la conservazione del riso, ancora in uso. Hanno la forma dell'isola di Lombok ed è diventata quindi una sorta di simbolo nazionale. Sono sollevate da terra con enormi pali di bambù per isolarle da topi, polli e uccelli.

Proseguiamo e scopriamo di essere particolarmente fortunati: essendo appena arrivati dei pullman di turisti, oggi si terrà una dimostrazione di antiche cerimonie locali a cui possiamo assistere liberamente. La prima, la Ghndà Belì, consiste in una danza effettuata da due giovani indigeni in abiti originali che suonano un enorme tamburo tradizionale, accompagnati da tre musicisti con un sottofondo tribale.

La seconda cerimonia è la Peresean Dance, e a quanto capisco dalla spiegazione è realizzata da "un referee and two painters" — un arbitro e due pittori. L'arbitro si presenta per primo accompagnato dalla musica, poi ai suoi lati compaiono due danzatori con due lunghi pennelli di bambù e due tele piatte legate al braccio. Ora, il loro inglese lascia a desiderare nella pronuncia, così non mi rendo conto che invece di "painters" aveva detto "fighters". Mi aspetto quindi che inizino a dipingere l'uno sulla tela dell'altro col lungo pennello, mentre invece improvvisamente utilizzano quei tubi di bambù per darsi mazzate con una veemenza incredibile: il pennello era in realtà un bastone e la tela uno scudo. Il suono dei colpi fa quasi paura, ma ancora più sconvolgente è che, terminata la danza, giunge il turno di due bambini che si esercitano a suon di mazzate.

Proseguiamo la visita verso i luoghi della comunità. Ci viene mostrata un'area rialzata coperta da un tetto di erba e bambù chiamata Bali Jajar — Bali Casa, Jajar Riunione — dove il capo del villaggio, eletto in maniera ereditaria, si ritrova a risolvere i problemi della comunità in qualità di mediatore. Un amministratore di condominio, insomma. Ad esempio, se due abitanti si sposano, coordina la comunità perché la nuova casa venga costruita da tutti insieme. La comunità conta 35 famiglie e 175 persone che vivono qui da cinque generazioni. I giovani oggi tendono a scegliere la modernità e una volta sposati se ne vanno, rendendo sempre più difficile mantenere la popolazione numerosa — si stima che senza questi spostamenti sarebbero oggi circa 500-600 persone. I legami familiari restano però saldi: chi parte è sempre il benvenuto a tornare per le feste e le cerimonie.

Assistiamo anche a una dimostrazione in cui una donna nativa spalma argilla e sterco di vacca sul pavimento: lo fanno circa due volte a settimana, lasciando seccare al sole e ripetendo l'operazione a strati. Come negli altri villaggi anche qui assistiamo alla creazione dei Sarong: ormai siamo diventati esperti, sia nella tecnica di tessitura che nello schivare i venditori di souvenir.

Parlando dei galli che ci tenevano svegli la notte, scopriamo che l'onomatopea "chicchirichì" per loro suona come "ciacciarló". Quando menzioniamo le moschee e i nostri timpani forati, Amat ci ride su e ci spiega che a Lombok ci sono circa mille moschee. Il 90% della popolazione è musulmana, mentre la parte restante è suddivisa tra cristiani e induisti che vivono principalmente a Mataram, per cui le aree rurali sono abitate quasi esclusivamente da musulmani.

Proseguiamo in una sorta di stalla dove riposano le mucche della comunità, usate per le cerimonie: quando c'è una festa o un matrimonio, ZAC! E tutti si mangia mucca. L'ultimo edificio che visitiamo è una torre di vedetta chiamata Gianga: durante la notte due vedette a turno scrutano il villaggio e in caso di problemi — un animale rubato, un incidente — suonano la campana e danno l'allarme.

Usciamo dal villaggio e, vinta dalla fame, Elisa si compra una pannocchia. Ormai non mi stupisce più niente. Ringraziamo la nostra guida, un altro esempio di forte accoglienza locale, e torniamo in macchina.

\begin{center}
    % Decorative line
    \noindent\rule{0.6\textwidth}{0.4pt}
\end{center}

Percorriamo un breve tratto di strada fino alla nostra ultima tappa: il villaggio di Sade. L'atmosfera qui è oggettivamente diversa e capiamo subito a cosa si riferiva Amat quando parlava di "mercatino a cielo aperto". L'accoglienza è sempre gentile ed educata, ma l'intento commerciale è evidente fin dall'ingresso: è richiesta una donazione all'entrata e una mancia alla guida al termine.

La nostra guida questa volta si chiama Aman, un abitante del villaggio un po' sopra le righe, con un inglese improvvisato di ardua interpretazione. Riporterò i fatti che ci ha raccontato così come li ha condivisi lui, e così come li ho capiti io, senza filtri.

Sade è il villaggio più antico di Lombok, fondato nel 1079, oggi alla quindicesima generazione, con 150 famiglie e 750 persone. Tutti i nativi discendono da un'unica famiglia: per mantenere la tradizione di sangue si sposano solo tra cugini. Sposarsi con esterni è possibile ma molto costoso — cinque mucche o cinque bufali. Se si sposano internamente non pagano, ma gli uomini spesso rapiscono le donne. Devo farglielo spiegare due volte perché non sono sicuro di aver capito bene: se il padre non acconsente a fare sposare la figlia, qualcuno la rapisce e la porta nella moschea dove lo sposo l'attende. Dopo il matrimonio segreto la coppia fugge lontano dal villaggio per non essere raggiunta dal padre. Una settimana dopo tornano e si celebra la cerimonia ufficiale con la famiglia. Ogni volta che rileggo questa parte penso sempre di aver capito male, ma ho registrato il racconto con il telefono. Aman ha effettivamente raccontato questa storia: il mio unico dubbio resta sulla sua conoscenza dell’inglese.

Iniziamo il nostro itinerario per le strette vie del villaggio, molto più densamente abitato rispetto ai precedenti e pieno di bancarelle, ma decisamente più povero. Presenta molti aspetti comuni agli altri: case coi tetti spioventi di erba che obbligano a inchinarsi per entrare, depositi di riso, torre di vedetta. E naturalmente le donne che tessono — anche qui Elisa proverà, aiutata dalla sorella di Aman.

Ormai sentiamo di sapere tutto sulla tessitura dei Sarong, così accettiamo volentieri la proposta di una tazza di caffè, per vivere un momento di comunità. Aman ritorna con due bicchieri di caffè fumante: è la prima volta che lo assaggiamo con lo zucchero di canna locale e ci pentiamo di non averlo fatto prima. Si tratta del solito caffè tostato col riso, ma la canna da zucchero lo rende dolce con una leggera sfumatura di terra e frutta secca. Mentre brindiamo ci insegnano il loro equivalente del "cin cin": nella lingua Sasak si dice "toss", molto simile al "make a toast" inglese.

Mentre beviamo e conversiamo, Aman si accende una sigaretta e gli chiedo di poter provare il loro tabacco. Gentilissimo, me ne rolla una tra le mani giunte in preghiera, senza filtro e senza bagnare la carta. La cartina è alla menta e insieme al tabacco locale dà alla sigaretta un gusto morbido e avvolgente.

Qui, riascoltando la registrazione, c'è uno scambio di battute che merita di essere tramandato.

— \textit{Why is the cigarette so sweet?} — chiedo alla guida. 

— \textit{Because the paper is sweet} — risponde Aman. 

— \textit{C'è il pepe?} — chiede Elisa.

Lo ammetto, il loro inglese talvolta è davvero difficile da comprendere.

Fumando tabacco e sorseggiando caffè ci sciogliamo e iniziamo a chiacchierare. Aman racconta di avere quattro figli, alcuni già sposati, e dei nipotini di quattordici, dodici e sei anni. Quando Elisa prova a fumare, notiamo una velata reazione di stupore da parte di Aman, così gli chiedo se qui le donne fumino. Ci comunica che in effetti non fumano, ma masticano le foglie di Betel.

Nota importante: inizialmente capiamo "foglie di tabacco" e solo riascoltando la registrazione comprendiamo quello che si rivelerà un errore irreparabile. Gli chiediamo subito se possiamo provarle anche noi e lui, come sempre gentilissimo, scompare per un attimo tornando con delle foglie verdi, una polvere bianca che ci comunica essere corallo bollito e tritato, e una noce di Areca, un nocciolo con la consistenza e la forma di uno spicchio d'aglio.

Elisa è rimasta particolarmente sconvolta dall'esperienza che segue, così la lascio raccontare a lei.

\begin{center}
    % Decorative line
    \noindent\rule{0.6\textwidth}{0.4pt}
\end{center}

\textit{"Il nocciolo di Areca si rivela commestibile e legnoso. Aman mi dice di metterlo in bocca e masticarlo. Essendo molto grande immagino intenda dirmi di assaggiarne un pezzo, ma sebbene glielo chieda più volte la risposta sembra essere sempre di mangiarlo tutto in un boccone.}

\textit{Appena inizio a masticarlo i miei occhi si sgranano e la salivazione aumenta in modo esponenziale: un gusto amaro che più amaro nella vita non si può mi pervade la bocca. Aman mi consegna le due foglie di Betel piegate con dentro la polvere di corallo e mi dice di mangiare anche quelle. Inizio a non sentire più la bocca ma non posso sputare perché temo di risultare maleducata e ingrata per il regalo, così continuo a masticare ingoiando tutto aiutandomi col caffè nero.}

\textit{el frattempo arriva un cane con il muso e la schiena completamente scuoiati e sanguinanti, affetto probabilmente da rabbia, peste, lebbra, scabbia, leishmaniosi o non so che altro, e sento gridare per la prima volta la sorella della guida, che allontana l'animale infetto: non sembrava nemmeno più lei, fino a un momento prima aveva una vocina delicata e soave. Mentre la povera bestia scappa io continuo ad avere una sensazione terribile in bocca e in tutto il corpo, sono stanca e spossata, seduta per terra sui gradini della casa, scalza, quando a un certo punto comprendo che è meglio proseguire.}

\begin{center}
    % Decorative line
    \noindent\rule{0.6\textwidth}{0.4pt}
\end{center}

Breve parentesi su cosa sia in realtà la noce di Areca Catechu, detta noce di Betel. 

\textit{La noce di Areca è il frutto della Areca catechu, che cresce in gran parte del Pacifico tropicale (Melanesia e Micronesia), sud-est asiatico, Asia meridionale, Asia orientale e parti dell'Africa orientale. Questo frutto è comunemente indicato come noce di Betel, quindi è facilmente confuso con foglie di Betel (Piper betle) che vengono spesso utilizzate per avvolgerlo (paan). Il termine Areca originato dalla parola Malayalam adakka e risale al XVI secolo, quando i marinai olandesi e portoghesi portarono la noce dal Kerala all'Europa.}

\textit{Il consumo ha molti effetti nocivi sulla salute ed è cancerogeno per l'uomo. Vari composti presenti nella noce, in particolare l'arecolina (il principale ingrediente psicoattivo che è simile alla nicotina), contribuiscono a cambiamenti istologici nella mucosa orale. Il consumo delle noci di Betel è una pratica molto antica. Notizie sugli usi e sulle proprietà dei semi di Betel sono state rinvenute in un documento cinese datato tra il 180 e il 140 a.C. La noce di Betel contiene molti tannini che macchiano i denti di nero. Queste sostanze hanno la proprietà di stimolare la secrezione salivare. Il modo più comune di consumare le noci di Betel è quello di tagliarle in FETTE SOTTILI, avvolgerle nelle foglie di pepe di Betel, preventivamente spolverate di calce, aggiungendo spezie come cannella, noce moscata, cardamomo, catecù perché DA SOLE SONO IMMANGIABILI. Masticare la noce di Betel, con o senza aggiunta di tabacco, provoca molteplici forme di cancro e malattie cardiovascolari. Il danno causato dal consumo di noci di Betel in tutto il mondo è stato classificato nel 2017 come "un'emergenza sanitaria pubblica globale trascurata".}

\begin{center}
    % Decorative line
    \noindent\rule{0.6\textwidth}{0.4pt}
\end{center}

Riprendo io la parola e ci tengo a precisare che Elisa è ancora viva, sta bene ed è qui a fianco a me sul divano mentre faccio la revisione del diario.

Il cane era davvero messo male, ormai più da abbattere che da curare. Le foglie di Betel da sole non erano male, ma ammetto di averle provate senza la noce. La guida è stata poi così gentile da girarci due sigarette pronte all'uso. Avevamo proposto di comprare il suo tabacco ma lo vendeva solo in confezioni da due etti, troppo per noi.

Anche la sorella è stata talmente disponibile con Elisa durante la tessitura che per poco non abbiamo comprato uno dei suoi lavori. Qui ammetto di aver visto Elisa fuori di sé — tra la noce di Betel e la sua febbre da souvenir non riusciva a tenersi a freno. Continuava a ripetere: "dobbiamo comprare qualcosa!"

No, no, fidati, non dobbiamo cara mia.

Dopo un momento di stasi in cui Elisa aveva indosso ben due abiti tradizionali pronti alla vendita, sono riuscito a contrattare di comprare solo un braccialetto: per me. La trattativa è cominciata da una cifra folle — capisco tenersi alti, ma sparare altissimo per poi scendere immediatamente a un quinto del prezzo mi sembra insensato. Oltretutto non avevano il resto, il che mi è sembrato un tentativo di farsi lasciare più soldi. Peccato, perché sono stati molto gentili a offrirci caffè, sigarette e foglie — o meglio, droghe leggere — ma se mi fossero venuti incontro con trattative più oneste sarei stato ben disposto a lasciare una mancia generosa. Non avendo il resto ho fatto un forfait e gli ho lasciato anche la mancia di fine tour.

Scampati al tentativo di scam e frenata Elisa che stava per acquistare un pacchetto di tre portachiavi con la scritta Lombok, proseguiamo fino alla torre di vedetta. La descrizione è la stessa del villaggio di Ende, ma qui ci è permesso di salire — non esattamente in tutta sicurezza, visto che la struttura scricchiola e cigola con tantissime riparazioni improvvisate. Dall'alto si gode di una spettacolare vista su tutto il villaggio e sulle aree limitrofe, che purtroppo Eli non si gode perché è ancora su di giri dalla noce di Betel. Scattiamo un po' di foto e penso che l'indomani, quando si sarà ripresa, dovrò raccontarle tutto: "guarda questa foto Eli, lo so che non ricordi ma eravamo quassù, in alto in alto."

Scendiamo tenendoci al corrimano di bambù che trema forse più degli scalini, e la tappa successiva è l'albero dell'amore, o "albero della fortuna", dove gli uomini fanno le promesse di matrimonio alle donne del villaggio. Deduco che sia solo per quelle che non sono state rapite.

Proseguendo entriamo in una capanna piena di Batik appesi al muro. Sono tutti bellissimi, ma accanto alle rappresentazioni originali ne troviamo molte di moderne e artefatte. Un nativo ci mostra la tecnica: prima realizzano uno schizzo su una tela di cotone, poi fanno bollire della cera naturale ottenendo una sostanza collosa che viene colorata con coloranti naturali e stesa sulla tela. La cera secca al sole, la parte liquida evapora e il colore resta impregnato nel cotone. Il processo si ripete più volte; alla fine la tela si mette a bollire per rimuovere tutta la cera rimasta. Il risultato è un tessuto utilizzabile da entrambi i lati, con il disegno che non perde definizione.

Non compriamo niente e proseguiamo fino all'uscita del villaggio, dove torreggia un'enorme scritta: "Matur Tampi Asih. Palah Bedait Malik" — "Thank you very much. See you next time." In un ultimo gesto di gentilezza, prima di lasciarci, Aman si mette in mezzo alla strada a bloccare auto e moto in arrivo per farci attraversare in sicurezza.

\begin{center}
    % Decorative line
    \noindent\rule{0.6\textwidth}{0.4pt}
\end{center}

Salutata la guida, notiamo che Abi è diventato più loquace. Proviamo a scambiare due parole ma a quanto intuiamo si è presumibilmente rotto le scatole. Senza dire molto ci accompagna al Novotel, lasciandomi ancora un briciolo di speranza per il coltello tradizionale e dicendomi che si informerà. Gli ricordo che non ha il mio numero di cellulare. Aspetto ancora fiducioso.

Arrivati al Novotel ci cambiamo e andiamo direttamente in spiaggia, dove ci corichiamo su due sdraio sotto un gazebo ispirato alle antiche abitazioni Sasak. Elisa dorme un po' e io inizio a scrivere il diario per fissare i ricordi infiniti della giornata. Sarebbe un bel momento di relax, se non fosse che vicino a noi c'è una coppia di calabresi che urla al vivavoce in telefonata con quello che è molto probabilmente un cugino che sta a Poggioreale. All'inizio penso stiano esagerando l'accento apposta, poi realizzo che parlano davvero come delle parodie. Questi potrebbero essere in Abruzzo e non se ne accorgerebbero nemmeno.

Durante la giornata abbiamo pranzato solo con una pannocchia bollita e delle patatine dolci piccanti comprate in una stazione di servizio, quindi optiamo per un aperitivo pre-cena. Al bar sulla spiaggia preferiamo però il venditore ambulante da cui acquistiamo acqua di cocco e un mango, che il ragazzo prepara di fronte a noi col machete. Ce li godiamo di fronte al tramonto.

Quando inizia a fare freddo dopo che il sole è sceso, Elisa mi chiede — puntuale — di fare il bagno. Non ne capisco il senso, ma poi realizzo: finché c'è il sole ci si abbronza, quando non c'è più si può fare altro. Non la seguo anche se mi giura che l'acqua è calda, e così rientra in modalità struzzo a cercare coralli sul fondale. Si rivela un'altra spedizione fallimentare come quella della tartaruga. Sta facendo buio quindi torniamo in camera a prepararci per uscire.

Questa sera vogliamo cenare a Kuta, lontano dall'hotel. Contatto Dian, l'autista che ci aveva accompagnati al Novotel dalla Gili Asahan e che nei giorni precedenti mi aveva scritto di poter contare su di lui per qualsiasi passaggio. È immediatamente disponibile e ci accordiamo sull'orario.

Arriviamo all'appuntamento in anticipo. Un autista ci fa segno di salire sulla sua auto e inizia a fare dei commenti che non comprendiamo bene. Inizia a sorgermi un dubbio e gli chiedo di ricordarmi il ristorante che mi aveva consigliato per messaggio: lui risponde prontamente, ma la risposta è completamente diversa, e capisco che siamo saliti sull'auto sbagliata. Ne ho la conferma quando arriviamo a Kuta e scopro che il suo numero di telefono è diverso da quello di Dian. E per finire l'autista si chiama Eric. Povero Dian! Appena scendiamo gli scrivo mortificato chiedendogli scusa, ma non mi risponderà mai più. Spero che la cosa si sia risolta con uno scambio e che anche lui abbia caricato i passeggeri che si erano accordati con Eric.

Come se non bastasse, durante il tragitto Eric ci aveva chiesto il numero di stanza al Novotel ed Elisa mi fa prendere male, sicura che l'abbia fatto per andare a rubarci in camera sapendo che siamo fuori a cena. Si tranquillizza solo quando le ricordo che tutto ciò che abbiamo di prezioso è con noi o nella cassaforte, ma ormai l'ansia si è trasferita addosso a me.

Giriamo un po' per Kuta che la sera ricorda molto Palma di Maiorca: piena di stranieri giovani attratti dalle spiagge e dal surf, che si mescolano con la tradizione locale. Non ci fa impazzire l'atmosfera, molto lontana dall'Indonesia che abbiamo visto finora — per non parlare dei villaggi visitati oggi, che hanno tutta un'altra sfumatura.

Dopo un paio di giri per la zona dei locali ci lasciamo attrarre da un ristorante di pesce: decisamente spartano e improvvisato, un gazebo con tavoli e sedie di bambù, griglia all'aperto e pesce fresco esposto lungo la strada. Un tavolino ricoperto di ghiaccio è pieno di pesci tropicali dai mille colori e dimensioni, tutti appena pescati da pescatori locali — non emanano alcun odore. Attendiamo un po' che si liberi un tavolo e alla fine riusciamo a sederci.

I prezzi del menù sono bassissimi, ma il pesce fresco va contrattato davanti alla piastra — cosa che purtroppo non ci viene in mente, trattandosi di un ristorante. Poco male: in qualunque altro posto del mondo il prezzo che otteniamo sarebbe stato bassissimo lo stesso, visto l'ordine di due aragoste, otto gamberoni, riso in bianco, insalata di verdure locali con glassa di soia e pomodori, birra e acqua. Il pesce è cucinato con sweet chili e pepe, risulta buonissimo e non dà fastidio al palato. Divoriamo tutto. Tra le recensioni molto positive che leggiamo sul posto ce n'è solo una con quattro stelle, dove ci si lamenta che la cucina è un po' "oleosa." In effetti mangiamo dentro ciotole di bambù intrecciato e l'olio del condimento passa attraverso la ciotola, la tovaglietta, il tavolo di bambù e mi cola sulle scarpe.

Ci alziamo soddisfatti e facciamo ancora una breve passeggiata prima di contattare Eric per venirci a prendere, questa volta memorizzando per bene il suo faccione dall'immagine del profilo WhatsApp. Arriva puntuale spaccando il minuto ed è effettivamente l'autista dell'andata. Come controprova non gli comunico l'hotel per vedere se sa dove andare. Ovviamente non se lo ricorda.

Arriviamo comunque al Novotel senza problemi e saliamo subito in camera a preparare le valigie per l'indomani, quando dovremo raggiungere le isole Gili — più precisamente Gili Trawangan, o Gili T — dove incontreremo Enrico e Cristina. Continuo a scrivere il diario finché ho i ricordi freschi e a un certo punto mi rendo conto che si è fatta l'una di notte. Riuscirò a dormire poco anche al Novotel.

\begin{center}
    % Decorative line
    \noindent\rule{0.6\textwidth}{0.4pt}
\end{center}

P.s. Una nota finale sull'esperienza di Elisa con la noce di Betel. Il metodo di ingerire la noce tagliata a fettine insieme alla foglia di Betel e alla polvere di corallo bollita è noto come "Betel Quid" ed è una pratica culturale spesso associata al consumo di stupefacenti. Elisa non l'ha presa a fette: ha mangiato una noce intera, associandola al tabacco e al caffè forte molto zuccherato. Ha vissuto un'esperienza mistica, profondamente amara. E io ora sono geloso.


\pagestyle{empty} % disable numbers

% da scegliere
\landscapeLargeMargin{images/day7/PXL_20240815_020256475.jpg}
%\landscapeLargeMargin{images/day7/PXL_20240815_024728464.jpg}
\landscapeLargeMargin{images/day7/PXL_20240815_024752346.MP.jpg}
\landscapeLargeMargin{images/day7/PXL_20240815_024843857.MP.jpg}
\landscapeLargeMargin{images/day7/PXL_20240815_053610118.MP.jpg}

%coppia
\twoImagesLargeMargins{images/day7/IMG_4538.JPG}
                      {images/day7/IMG_4537.JPG}
\landscapeLargeMargin{images/day7/IMG_4583.JPG}
%

%coppia
\twoImagesLargeMargins{images/day7/IMG_4622.JPG}
                      {images/day7/IMG_4626.JPG}
\twoImagesLargeMargins{images/day7/IMG_4738.JPG}
                      {images/day7/IMG_4732.JPG}
%

%coppia
\twoImagesLargeMargins{images/day7/IMG_4698.JPG}
                      {images/day7/IMG_4696.JPG}
\landscapeThinMarginCentered{images/day7/IMG_4702.JPG}
%

%coppia
\threeImagesLargeMargins{images/day7/PXL_20240815_053431365.MP.jpg}
                        {images/day7/PXL_20240815_053343669.MP.jpg}
                        {images/day7/PXL_20240815_053354968.jpg}
\landscapeLargeMargin{images/day7/PXL_20240815_064120730.MP.jpg}
%

%coppia
\splitImageLeft{images/day7/IMG_4903.JPG}
\splitImageRight{images/day7/IMG_4903.JPG}
%

%coppia
\portraitFullscreen{images/day7/IMG_4908.JPG}
\landscapeThinMarginCentered{images/day7/IMG_5000.JPG}


%%% TODO: ripartire da qui
%coppia
\landscapeThinMarginCentered{images/day7/IMG_4987.JPG}
\landscapeThinMarginCentered{images/day7/IMG_4987.JPG}
%

%coppia
\fourPortraitImages{images/day7/PXL_20240815_064717667.jpg}
                   {images/day7/PXL_20240815_065219166.jpg}
                   {images/day7/PXL_20240815_065233804.MP.jpg}
                   {images/day7/PXL_20240815_065256568.jpg}
\twoImagesLargeMargins{images/day7/PXL_20240815_072317333.MP.jpg}
                      {images/day7/PXL_20240815_072331184.MP.jpg}
%

\twoImagesLargeMargins{images/day7/PXL_20240815_072936448.jpg}
                      {images/day7/PXL_20240815_072845411.MP.jpg}
\landscapeLargeMargin{images/day7/PXL_20240815_073327079.MP.jpg}


%coppia
\splitImageLeft{images/day7/PXL_20240815_101911184.jpg}
\splitImageRight{images/day7/PXL_20240815_101911184.jpg}
%

%coppia
\landscapeLargeMarginCentered{images/day7/PXL_20240815_095152760.PORTRAIT.jpg}
\landscapeLargeMarginCentered{images/day7/PXL_20240815_100233901.PORTRAIT.jpg}
%

%coppia
\landscapeThinMargin{images/day7/PXL_20240815_102344296.MP.jpg}
\twoImagesLargeMargins{images/day7/PXL_20240815_101305474.MP.jpg}
                      {images/day7/PXL_20240815_102410573.MP.jpg}
%

%coppia
\landscapeLargeMargin{images/day7/PXL_20240815_124156382.MP.jpg}
%

\threeImagesLargeMargins{images/day7/PXL_20240815_130610813.jpg}
                        {images/day7/PXL_20240815_130550458.MP.jpg}
                        {images/day7/PXL_20240815_130813894.jpg}
%

%coppia
\landscapeThinMargin{images/day7/PXL_20240815_115458733.MP.jpg}
\twoImagesLargeMargins{images/day7/PXL_20240815_133333056.jpg}
                      {images/day7/PXL_20240815_133359565.jpg}
%

\landscapeThinMarginCentered{images/day7/PXL_20240815_133533090.MP.jpg}

\pagestyle{fancy} % enable numbers